\section{Introduction}
\label{sec:introduction}

Mobile agents increasingly operate applications by reading screenshots or accessibility trees and producing taps, swipes, and text-entry actions. Systems and benchmarks such as AppAgent, Mobile-Agent, AutoDroid, AndroidWorld, AndroidEnv, Android in the Wild, and SPA-Bench show why this interface is attractive: it reaches applications without requiring each developer to expose an agent-specific API \cite{zhang2023appagent,wang2024mobileagent,wen2023autodroid,rawles2024androidworld,toyama2021androidenv,rawles2023aitw,chen2024spabench}. The same interface collapses a security boundary. If a code or credential is visible enough for an agent to copy, it is visible enough for a compromised model, cloud endpoint, injected webpage, malicious notification, or logging layer to exfiltrate.

The risk is not hypothetical at the level of mechanism. Indirect prompt injection lets untrusted content compete with the user's instruction \cite{greshake2023indirect,zhan2024injecagent,debenedetti2024agentdojo}; GUI-specific studies demonstrate that visually embedded or third-party content can steer agents toward privacy violations \cite{chen2025fineprint,liao2024eia,chen2025aeia}. A recent mobile study exercises multi-app adversarial workflows and cross-app data exfiltration \cite{du2025mobileattacks}. These results do not imply that every agent or every task fails, but they establish that a secret exposed to the planning loop is inside the attack surface.

Prompt filtering cannot provide the required invariant. A detector may miss a sensitive span; a placeholder may be substituted into the wrong field; and a model that never reads a value may still direct a local resolver to disclose it at an attacker-chosen destination. Mobile placeholder privacy is an important step: the closest design transforms sensitive screen content into type-preserving placeholders and mediates actions locally \cite{zhao2026anonymization}. AutoDroid also includes prompt-side substitution as part of mobile task automation \cite{wen2023autodroid}. Both motivate an abstraction boundary, but a security claim that begins after screen capture depends on detection coverage and does not itself bind the eventual value to the app, web origin, task, or field that is authorized to receive it.

This paper asks a narrower systems question: \emph{can Android let an untrusted agent use a secret without ever giving the agent the secret or a transferable bearer copy?} Our answer is \SH, an OS-minted opaque capability. Android already contains authoritative source points and relevant enforcement primitives: SMS Retriever routes app-associated verification messages \cite{androidSmsRetriever}; Autofill and Credential Manager mediate user data for app fields \cite{androidAutofill,androidCredentialManager,androidAutofillServices}; Android 15 redacts detected OTP content from untrusted notification listeners and remote screen sharing \cite{android15privacy}; and hardware-backed Keystore, key attestation, and Protected Confirmation support protected keys and high-assurance confirmation \cite{androidKeystore,androidKeyAttestation,androidProtectedConfirmation}. \SH\ connects these ideas to the accessibility-agent action path.

At mint time, the trusted source pipeline gives the SecretHandle service a raw value, a type, provenance, and the narrowest defensible destination. The service stores the value in OS-private memory or sealed storage and returns a randomized identifier plus an authenticator. The agent receives only a typed description such as \texttt{OTP\#7 expires in 43s; destination: com.example.bank/login\_otp}. To act, it sends \texttt{ENTER(OTP\#7)} and a target-node token. The resolver authenticates the handle, verifies the current task, package, WebView origin when trustworthy, field class, time, remaining uses, and confirmation rule, atomically reserves the record, and injects the value below the agent's observation boundary. The subsequent accessibility tree and screenshot contain a masked projection rather than the value.

Destination binding resembles audience restriction, not encryption for the model. OAuth resource indicators prevent a token intended for one resource from being replayed at another \cite{rfc8707}; DPoP additionally constrains a token to a presenter key \cite{rfc9449}. WebAuthn scopes credentials to relying parties \cite{webauthn3}. \SH\ applies a related principle to local UI data flow: possession of a handle is insufficient unless the OS-observed target and task match the handle policy. This analogy guides the design but does not make an OAuth token, passkey, or payment token an empirical baseline for Android text entry.

\subsection{Contributions}

This manuscript makes five scoped contributions.

\begin{enumerate}[leftmargin=*]
  \item It defines an OS-minted secret-handle abstraction for Android in which the raw value is intercepted at an authoritative source and represented to a mobile AI agent by a non-semantic, authenticated capability.
  \item It designs a destination- and task-bound resolver for \texttt{ENTER(handle)} with atomic use reservation, typed lifecycles, optional protected confirmation, and explicit failure semantics. The destination includes package identity and may include a trustworthy web origin and field class; this granularity is a policy decision, not assumed ground truth.
  \item It integrates observation-side redaction across accessibility and screenshot channels and identifies the closure conditions required to keep a just-entered value out of the next planning observation.
  \item It supplies a formal transition system and conditional proofs for non-disclosure in the modeled agent view, destination confinement, reduction-style unforgeability, and single-use/expiry safety. The results do not claim to stop malicious actions that remain within an authorized destination.
  \item It gives an implementation blueprint and preregistration-style evaluation plan for AOSP devices, login/2FA/checkout applications, real mobile agents, redirection and replay attacks, ablations, leakage instrumentation, and latency measurement. No unmeasured number is presented as a result.
\end{enumerate}

\subsection{Scope and paper organization}

The trusted computing base includes the Android kernel and relevant framework services, the source adapters, the SecretHandle store and resolver, the redaction paths, and protected key operations. The agent, model provider, third-party content, and ordinary applications are untrusted. A compromised OS, a covert camera, timing leakage, or a malicious authorized destination is outside the central guarantee and is discussed in \cref{sec:limitations}. The rest of the paper separates prior work (\cref{sec:related}), notation and platform context (\cref{sec:preliminaries}), the model and properties (\cref{sec:system,sec:problem}), the design and algorithms (\cref{sec:method,sec:algorithm}), formal analysis (\cref{sec:theory}), and the planned evaluation (\cref{sec:experiments}).
